%!TeX program = xelatex
% =============================================================================
% BEAMER THEME CONFIGURATION
% Last Modified: 2026
% Description: Professional presentation template with custom TikZ headers,
%              footers, and auto-scaling components.
% =============================================================================

\documentclass[aspectratio=169]{beamer}

% -----------------------------------------------------------------------------
% Core Packages
% -----------------------------------------------------------------------------
\usepackage{amsmath, amssymb, xcolor, tikz}
\usetikzlibrary{calc, backgrounds}
\usepackage[most]{tcolorbox}
\tcbuselibrary{skins}           % Required for 'enhanced' tcolorbox styles
\usepackage{xparse}             % Advanced command definitions
\usepackage{fontspec}           % Font selection for XeLaTeX
\usepackage{xeCJK}              % Chinese language support
\usepackage{unicode-math}       % Unicode mathematics support
\usepackage{etoolbox}           % Logical toggles and patching
\usepackage{booktabs}           % Professional table horizontal rules
\usepackage{pgfplots}           % Vector graphics for plots
\usepackage{adjustbox}          % Content scaling and overflow prevention

\pgfplotsset{compat=1.18}

% -----------------------------------------------------------------------------
% UI/UX Settings
% -----------------------------------------------------------------------------
% Disable standard Beamer navigation symbols for a cleaner interface
\setbeamertemplate{navigation symbols}{}

% Define Custom Color Palette
\definecolor{lightblue}{RGB}{70,130,180}
\definecolor{midblue}{RGB}{28,70,130}
\definecolor{darkblue}{RGB}{10,40,80}
\definecolor{myred}{RGB}{150,50,30}
\definecolor{mygreen}{RGB}{45,125,77}
\definecolor{mypurple}{RGB}{120,61,137}
\definecolor{myorange}{RGB}{190,88,32}
\definecolor{myitem}{RGB}{35,35,35}

% -----------------------------------------------------------------------------
% Typography Configuration (XeLaTeX)
% -----------------------------------------------------------------------------
\newtoggle{fontBlackTheme}
\AtBeginDocument{
  \iftoggle{fontBlackTheme}{
    % ========== 黑体主题 ==========
    \IfFontExistsTF{Microsoft YaHei}{
        % Windows
        \setCJKmainfont{Microsoft YaHei}
        \setCJKsansfont{Microsoft YaHei}
        \setCJKmonofont{Microsoft YaHei}
    }{
        \IfFontExistsTF{STHeiti-SC}{
            % macOS
            \setCJKmainfont{STHeiti-SC}
            \setCJKsansfont{STHeiti-SC}
            \setCJKmonofont{STHeiti-SC}
        }{
            \IfFontExistsTF{Noto Sans CJK SC}{
                % Linux / CI
                \setCJKmainfont{Noto Sans CJK SC}[AutoFakeBold=true]
                \setCJKsansfont{Noto Sans CJK SC}[AutoFakeBold=true]
                \setCJKmonofont{Noto Sans CJK SC}
            }{
                \IfFontExistsTF{FandolHei}{
                    % Linux fallback
                    \setCJKmainfont{FandolHei}[AutoFakeBold=true]
                    \setCJKsansfont{FandolHei}[AutoFakeBold=true]
                    \setCJKmonofont{FandolHei}
                }
            }
        }
    }
  }{
    % ========== 宋体主题 ==========
    \IfFontExistsTF{SimSun}{
        % Windows
        \setCJKmainfont{SimSun}[AutoFakeBold=true]
        \setCJKsansfont{SimSun}[AutoFakeBold=true]
    }{
        \IfFontExistsTF{STSongti-SC}{
            % macOS
            \setCJKmainfont{STSongti-SC}[AutoFakeBold=true]
            \setCJKsansfont{STHeiti-SC}[AutoFakeBold=true]
        }{
            \IfFontExistsTF{Noto Serif CJK SC}{
                % Linux / CI
                \setCJKmainfont{Noto Serif CJK SC}[AutoFakeBold=true]
                \setCJKsansfont{Noto Sans CJK SC}[AutoFakeBold=true]
            }{
                \IfFontExistsTF{FandolSong}{
                    % Linux fallback
                    \setCJKmainfont{FandolSong}[AutoFakeBold=true]
                    \setCJKsansfont{FandolSong}[AutoFakeBold=true]
                }
            }
        }
    }
  }
}
\newcommand{\UseBlackFont}{\toggletrue{fontBlackTheme}}
\newcommand{\UseSongFont}{\togglefalse{fontBlackTheme}}


% -----------------------------------------------------------------------------
% Layout Templates (Header & Footer)
% -----------------------------------------------------------------------------
% Custom Headline: Provides vertical padding
\setbeamertemplate{headline}{\vspace{0.8cm}}

% Custom Frame Title: Renders a TikZ-based header with a divider line
\setbeamertemplate{frametitle}{
    \begin{tikzpicture}[remember picture,overlay]
        \clip (current page.north west) rectangle ([yshift=-1.5cm]current page.north east);
        \node[anchor=west, font=\normalsize\bfseries, text=black, xshift=0.5cm, yshift=-0.5cm]
              at (current page.north west){\insertframetitle};
        \draw[line width=1pt, color=midblue]
              ([yshift=-0.85cm, xshift=0.3cm]current page.north west) -- ([yshift=-0.85cm, xshift=-0.3cm]current page.north east);
    \end{tikzpicture}
}

% Custom Footline: Tricolor block design with project meta-data
\setbeamertemplate{footline}{%
  \ifnum\value{framenumber}>0
    \begin{beamercolorbox}[wd=\paperwidth,ht=0.5cm,dp=0ex]{footline}
      \begin{tikzpicture}[remember picture,overlay]
        % Background Blocks
        \node[fill=lightblue, minimum width=0.3\paperwidth, minimum height=0.5cm, anchor=south west]
             at (current page.south west) (leftblock) {};
        \node[fill=midblue, minimum width=0.6\paperwidth, minimum height=0.5cm, anchor=south west]
             at ([xshift=0.3\paperwidth]current page.south west) (midblock) {};
        \node[fill=darkblue, minimum width=0.1\paperwidth, minimum height=0.5cm, anchor=south west]
             at ([xshift=0.9\paperwidth]current page.south west) (rightblock) {};

        % Text Labels
        \node[anchor=center, font=\tiny, text=white] at (leftblock.center) {PROJECT NAME};
        \node[anchor=center, font=\tiny, text=white] at (midblock.center) {INTERNAL PRESENTATION 2026};
        \node[anchor=center, font=\tiny, text=white] at (rightblock.center) {\insertframenumber};
      \end{tikzpicture}
    \end{beamercolorbox}
  \fi
}

% -----------------------------------------------------------------------------
% Custom Components & Environments
% -----------------------------------------------------------------------------
% Itemize Style: Small circular bullets using TikZ
\newcommand{\useSmallCircleItem}{%
  \setbeamertemplate{itemize item}{\raisebox{0.35ex}{\tikz{\draw[myitem, fill=myitem, opacity=0.8] (0,0) circle (1pt);}}}
  \setbeamertemplate{itemize subitem}{\raisebox{0.35ex}{\tikz{\draw[myitem, fill=myitem, opacity=0.8] (0,0) circle (1pt);}}}
}

% Block Styling: Custom tcolorbox wrappers for standard Beamer blocks
\newtcolorbox{beamerblockbox}[3][]{%
  enhanced, colback=#2!10, colframe=#2, coltitle=white, title=#3,
  fonttitle=\bfseries, arc=1mm, boxrule=0.5pt, left=2mm, right=2mm, top=1mm, bottom=1mm, #1
}
\renewenvironment{block}[1]{\begin{beamerblockbox}{midblue}{#1}}{\end{beamerblockbox}}
\renewenvironment{alertblock}[1]{\begin{beamerblockbox}{myred}{#1}}{\end{beamerblockbox}}
\renewenvironment{exampleblock}[1]{\begin{beamerblockbox}{myorange}{#1}}{\end{beamerblockbox}}

% Table Scaling: Wrapper to prevent table overflow
\newcommand{\ThreeLineTable}[1]{%
  \begin{adjustbox}{max width=\textwidth}{\footnotesize #1}\end{adjustbox}%
}

\useSmallCircleItem
\UseSongFont % Utilizing Serif/Songti theme for a professional look

\begin{document}

% --- Title Page ---
% Replace these with your actual details
\newcommand{\mytitle}{Investigation of [Your Subject] Using [Your Method]}
\newcommand{\myauthor}{xxxxxxxxxxxxxxxx}
\newcommand{\myinstitute}{xxxxxxxxxxxx \\ xxxxxxxxxxxxxxxxx}

\begin{frame}[plain]
  \begin{tikzpicture}[remember picture, overlay]
    % Header Decoration
    \node[fill=midblue, minimum width=\paperwidth, minimum height=0.5\paperheight, anchor=north]
          at ([yshift=+1pt]current page.north) {};

    % Main Title
    \node[text=white, font=\fontsize{28}{32}\selectfont\bfseries, text width=0.8\paperwidth, align=center]
          at ([yshift=-0.25\paperheight]current page.north) {\mytitle};

    % Credentials
    \node[text=black, font=\Large\bfseries] at ([yshift=-1.5cm]current page.center) {\myauthor};
    \node[text=black, font=\large, align=center] at ([yshift=-2.8cm]current page.center) {\myinstitute};
  \end{tikzpicture}
\end{frame}

% --- Table of Contents ---
\begin{frame}{Presentation Roadmap}
  \vspace{1em}
  \begin{enumerate}
    \item zzzzzzzzzzzzzzzz
    \item xxxxxxxxxxxxxxxx
    \item yyyyyyyyyyyyyyyy
    \item oooooooooooooooo
  \end{enumerate}
\end{frame}

% --- Equation Derivation Page ---
\begin{frame}{Performance Assessment Metrics}
  \textbf{Global Efficiency Score}: The following metric accounts for both modular performance and stochastic environmental variance.

  \[
    \boxed{
      S = \underbrace{\sum_{i=1}^{N} \textcolor{myorange}{\omega_i}\, \textcolor{mygreen}{\sigma_i}}_{\text{Weighted Modular Score}} \;+\;
      \overbrace{\left( {\frac{1}{N} \sum_{j=1}^{M} \textcolor{myred}{\delta_j} \cdot \prod_{k=1}^{K} \phi(x_k)} \right)}^{\text{Environmental Correction Factor}}
    }
  \]

  \[
    \left\{
    \begin{aligned}
      \textcolor{mygreen}{\sigma_i} &\;=\; \text{Performance coefficient for the } i\text{-th localized module}, \\
      \textcolor{myorange}{\omega_i} &\;=\; \text{Dynamic weight optimized via gradient descent}, \\
      \textcolor{myred}{\delta_j} &\;=\; \text{Error residual attributed to environmental factor } j.
    \end{aligned}
    \right.
  \]

  \vspace{0.5em}
  \small \textbf{Note:} The correction factor is auto-scaled to maintain legibility regardless of term length.
\end{frame}

% --- Block Emphasis Page ---
\begin{frame}{Optimization Strategies}
  Standardization of emphasis components ensures visual consistency across the deck.

  \begin{block}{Model Convergence}
    We assume the loss function is Lipschitz continuous to ensure global convergence under the specified learning rate.
  \end{block}

  \begin{alertblock}{Critical Vulnerability}
    Vanishing gradients may occur if the activation functions are not properly normalized during the initial heating phase.
  \end{alertblock}

  \begin{exampleblock}{Key Technical Insights}
    \begin{itemize}
        \item \textbf{Consistency:} Visual blocks allow for rapid scanning of core conclusions.
        \item \textbf{Safety:} Alert blocks should be reserved for edge cases and safety constraints.
    \end{itemize}
  \end{exampleblock}
\end{frame}

% --- Q&A Frame ---
\begin{frame}{Discussion}
  \centering
  \vspace{0.8cm}
  {\Huge \bfseries Q\&A}
  \vspace{1cm}

  \textcolor{midblue}{Thank you for your attention.} \\
  \texttt{xxxxxxxxxxxxxxx.org}
\end{frame}

% --- Ending Frame ---
\begin{frame}[plain]
  \begin{tikzpicture}[remember picture,overlay]
    \node[fill=midblue, minimum width=\paperwidth, minimum height=\paperheight] at (current page.center) {};
    \node[font=\fontsize{40}{70}\selectfont\bfseries, text=white] at (current page.center) {FIN};
  \end{tikzpicture}
\end{frame}

\end{document}